\documentclass[10pt]{article}
\usepackage[margin=0.78in]{geometry}
\usepackage{booktabs}
\usepackage{graphicx}
\usepackage{hyperref}
\usepackage{xcolor}
\usepackage{float}
\usepackage{caption}
\usepackage{array}
\hypersetup{colorlinks=true,linkcolor=blue!55!black,citecolor=blue!55!black,urlcolor=blue!55!black}
\input{generated/macros.tex}
\title{Open, Limp or Fold at 20 Big Blinds:\\
Testing Hwang's PLO Starting-Hand Tiers and Proposing Solver-Aligned Archetypes}
\author{Nuttakit Kundum}
\date{September 2026}

\begin{document}
\maketitle

\begin{abstract}
Jeff Hwang's Premium, Speculative, Marginal and Trash tiers are a common way to choose starting
hands in pot-limit Omaha. We ask how well they describe what a solver does with the first voluntary
action of a six-handed hand at 20 big blinds, when the pot is unopened and the player can fold,
limp or make a pot-sized raise. We solve a game with pot-sized betting on every street and all-in
bounded re-raising, using an abstraction that never places two Hwang tiers in one bucket, and we
audit exact hands by comparing the value of folding, limping and raising. Three findings follow.
First, the tiers order hands correctly but do not act as position-free rules: \PremFoldUTG\% of
Premium combinations are proven folds under the gun, all of them low straight hands that are not
double-suited. Second, the first action is driven mainly by the number of cards Ten or higher,
then by whether a suited Ace is present, then by pair rank; straight connectivity, the centre of
Hwang's scheme, has little independent effect at this depth. Third, \ArchCount\ archetypes built
from those drivers reach $R^2 = \RsqArchUTG$ for the solver's under-the-gun fold rate, against
$R^2 = \RsqTierUTG$ for Hwang's tiers, and cut the EV lost to wrong first-in decisions from
\LossHwangUTG\ to \LossArchUTG\ big blinds per decision on held-out hands.
\end{abstract}

\section{Question}
Hwang sorts playable four-card hands by structure: big cards and Ace-high Broadway wraps,
straight hands, suited-Ace hands, pair-plus hands, Aces, and marginal hands \cite{hwang}. Each
form is then graded Premium, Speculative, Marginal or Trash. The project's classifier implements
that grading from ranks, gaps and suits. It labels 5.49\% of the 270{,}725 physical hands Premium,
13.83\% Speculative, 19.91\% Marginal and 60.76\% Trash.

We started from a narrow claim: a Premium hand should never be folded when the pot is unopened and
the player is first to act. We widen it to three questions. Which hands does a solver open, limp or
fold at each seat? Where does Hwang's grading disagree with it? Can a small set of new archetypes
describe the solver better?

Hwang's own advice is conditional. He writes that ``the premium-class hands can usually be played
from any position and under most circumstances,'' that one should ``usually limp from up front,''
and that Marginal hands are ``late-position, minimum bet only'' \cite{hwang}. His book targets deep,
multi-way cash games. A 20-big-blind stack is a different setting, and part of the point is to see
which of his ideas survive it.

\section{Game and solver}
Six players, 0.5/1 blinds, 20 big blinds each, no ante, no rake, chip-EV payoffs. Every raise is
pot-sized and clipped to all-in. Preflop allows three raises (open, three-bet and a four-bet that is
effectively a shove). The flop, turn and river each allow a pot-sized bet and two raises. Side pots,
split pots and Omaha's two-plus-three showdown rule are exact. The public tree has 474{,}658 nodes
and 232{,}258 decisions.

\paragraph{Abstraction.} Preflop, 572 structural buckets are split by Hwang tier into 780 buckets,
so a bucket never mixes tiers. On the flop and turn, hands are grouped by the share of two-card
holdings that beat them on the current board (10 bins, with a separate nut bin), by flush draw
(none, non-nut, nut) and by the number of ranks that complete a straight on the next card
(0, 1, 2, 3+), giving 120 buckets per street; the river uses strength alone.

\paragraph{Training.} External-sampling Monte Carlo CFR with linear discounting, run as lock-free
parallel updates on a 14-core Apple M4 Pro. Two independent seeds were trained for 100 minutes
each, \Deals\ million deals per seed with six traversals per deal. The two seeds' tier-level
first-in frequencies agree to within 3.1 points at every seat except the small blind.

\paragraph{Validation.} The Omaha evaluator agrees with \texttt{phevaluator}; buckets are checked to
be tier-pure; chips are conserved at every terminal; closed-form cases (everyone folds to an open:
$+1.5$ big blinds) are exact. For 150 random hands per seat, choosing a better first action gains
at most 0.05 big blinds on average, so the first-in decisions are close to self-consistent. This is
a lower bound on exploitability, not a certificate.

\section{Two kinds of evidence}
\paragraph{Solver frequencies.} For every one of the 16{,}432 suit-isomorphism classes we record the
solved fold, limp and raise probabilities at each first-in seat, averaged over both seeds. These
cover every hand exactly, but the solver chooses per bucket, so two hands in one bucket always get
the same mix.

\paragraph{Per-hand audit.} For an exact hand at its first-in node we estimate the value of folding,
limping and raising when every player, the hand included, follows the solved strategy afterwards;
deals are weighted by the probability that every earlier seat folded. A hand is \emph{enter
proven} if limping or raising beats folding by at least three standard errors, \emph{fold proven}
if both lose to folding by at least three standard errors, and \emph{too close} otherwise. All 922
Premium classes and 400 uniformly drawn hands from each other tier were audited at all five seats:
\AuditRows\ hand--seat rows over \AuditClasses\ classes, weighted to the full population. Across the
two seeds, \SeedContra\ of \SeedPairs\ Premium hands under the gun received opposite proven verdicts.
Because the audit is per hand, it can separate hands that share a bucket: KK32 single-suited gains
$+0.30$ big blinds by limping under the gun while TT32 loses $0.93$, although the solver plays both
the same way.

\section{How often each tier folds}
\begin{table}[H]
\centering
\caption{Solved first-in mix, fold/limp/raise (\%), mean of two seeds.}
\label{tab:tiers}
\small
\input{generated/table_tiers.tex}
\end{table}

\begin{table}[H]
\centering
\caption{Per-hand audit: share of combinations proven to enter / too close / proven to fold (\%).}
\label{tab:verdicts}
\small
\input{generated/table_verdicts.tex}
\end{table}

\begin{figure}[H]
\centering
\includegraphics[width=\linewidth]{generated/fig_verdicts.pdf}
\caption{Per-hand verdicts by tier and seat.}
\label{fig:verdicts}
\end{figure}

The tiers keep their order at every seat. Premium is folded least and Trash most, and all tiers are
played more as the seat moves toward the button. Under the gun, \PremEnterUTG\% of Premium combinations
must enter and \PremFoldUTG\% must fold; on the button the proven-fold share falls to \PremFoldBTN\%.
Trash is a proven fold for \TrashFoldUTG\% of combinations under the gun. The small blind is
different for every tier: facing only the big blind with half a big blind already posted, the
solver limps a large share of all hands, including 45\% of Trash, and very few hands are proven folds.
Most Marginal hands, and many Speculative ones, are too close to call: at 20 big blinds they are near
break-even decisions rather than clear folds.

\section{What decides open, limp or fold}
Table~\ref{tab:drivers} splits all hands by one feature at a time. ``Play'' is the solver's
probability of limping or raising; ``entry EV'' is the audited value of the better of limping and
raising, minus folding, in big blinds.

\begin{table}[H]
\centering
\caption{First-in play rate and entry value by single features.}
\label{tab:drivers}
\small
\input{generated/table_drivers.tex}
\end{table}

\paragraph{Big cards come first.} Hands with zero or one card Ten or higher are almost never played
from early position and lose close to a big blind when they are; with three big cards the solver
enters 69\% under the gun and 96\% on the button; with four, nearly always. No other single feature
separates hands as sharply.

\paragraph{The Ace suit matters more than the second suit.} A suited Ace lifts under-the-gun entry
value from $-0.45$ (best suit Ten to King) to $+0.24$. A second suit helps (double-suited $+0.12$
against single-suited $-0.53$), and a hand whose best suit is headed by a Nine or lower plays almost
like a rainbow hand.

\paragraph{Pairs follow rank.} Aces are always played and are worth about three big blinds first in.
Kings and Queens are profitable from every seat. Jacks and Tens are near break-even under the gun
($-0.12$) and profitable only later. Small pairs are folds except in the small blind.

\paragraph{Connectivity matters little on its own.} Four ranks inside a five-rank window raise the
under-the-gun play rate only to 35\%, and entry value stays negative. Rundowns are valuable when they
also carry big cards and suits, which is the interaction our archetypes use.

\paragraph{Limp or raise.} Aces and hands with three or more big cards and a suited Ace or two
suits are raised far more often than limped at every seat, and raising is worth more than limping
by about 0.15--0.3 big blinds. Big pairs without a suited Ace are limped about 60\% of the time under
the gun, although the audited values of limping and raising differ by only about 0.1 big blinds.
Broadway rundowns are the one group where limping under the gun is clearly worth more
($+0.22$ against $+0.12$), which fits Hwang's advice to keep drawing hands multi-way from early
position. From the small blind, hands that enter are raised more often than limped, except the
weakest groups.

\section{Where Hwang disagrees with the solver}
Table~\ref{tab:mismatch} lists the largest groups (Hwang tier, form and suit pattern) where Hwang's
advice and the solver point in opposite directions under the gun. We read Hwang's advice as: raise
Premium, limp Speculative, play Marginal only from the cutoff, button and small blind, and fold Trash.

\begin{table}[H]
\centering
\caption{Largest disagreements under the gun. Entry EV is the audited mean where available.}
\label{tab:mismatch}
\small
\input{generated/table_mismatch.tex}
\end{table}

Three patterns explain most of the disagreement.
\begin{itemize}
\item \textbf{Straight and pair-plus-connector hands are overrated.} Single-suited pair-plus-connector
hands (for example 9-9-8-7) and gapped straight hands without big cards are Speculative or Premium in
Hwang's scheme, but under the gun they lose about 0.35--0.7 big blinds against folding. These are the forms
that rely on deep stacks and implied odds.
\item \textbf{Big-card hands with a dangler are underrated.} Big pairs with danglers, three Broadway
cards with a dangler, and weak suited-Ace hands that are double-suited are Marginal for Hwang, yet the
solver enters most of them under the gun, and the double-suited ones are profitable
($+0.17$ to $+0.66$). At 20 big blinds high-card strength is often realised by showdown or by a
single flop commitment, so a dangler costs less than it does deep.
\item \textbf{Suit quality is too coarse.} The classifier treats any two cards of a suit as
``suited''. The solver separates a suited Ace, a suited King to Ten, and low suits sharply, and it
plays three-flush hands almost as rarely as rainbow ones (9\% and 6\% under the gun, against 20\%
for single-suited hands).
\end{itemize}

\begin{figure}[H]
\centering
\includegraphics[width=0.52\linewidth]{generated/fig_rundowns.pdf}
\caption{Premium straight hands under the gun: best entry minus fold, in big blinds, by rank set and
suit pattern (ds double-suited, ss single-suited, 3f three-flush, mono monotone).}
\label{fig:rundowns}
\end{figure}

Figure~\ref{fig:rundowns} shows the Premium case in detail. Double-suited rundowns are near break-even
or better, and rundowns headed by two big cards (J-T-9-8, Q-J-T-9, K-Q-J-9) are playable when suited.
Single-suited, three-flush and monotone rundowns from 6-5-3-2 to 9-8-7-6 lose 0.4--1.3 big blinds
against folding. These hands carry the Premium label because the classifier grades straight forms by
their gaps and treats any suit as enough, while at this depth their equity when called is too low.

\section{Solver-aligned archetypes}
The drivers suggest grouping hands by big cards, Ace suit, second suit and pair rank, with
connectivity used only in combination with big cards. Table~\ref{tab:rules} gives the rules, checked
in order; Table~\ref{tab:archetypes} gives the recommended first action at each seat, taken from the
audited mean values of limping and raising. A cell is ``mixed'' when the better entry is within 0.10
big blinds of folding.

\begin{table}[H]
\centering
\caption{Archetype rules. A big card is a Ten or higher; a big pair counts as two big cards; a
rundown has all four ranks inside a five-rank window.}
\label{tab:rules}
\small
\input{generated/table_archetype_rules.tex}
\end{table}

\begin{table}[H]
\centering
\caption{Archetypes: share of hands, Hwang tiers they contain (P, S, M, T in \%), and the recommended
first action by seat.}
\label{tab:archetypes}
\small
\input{generated/table_archetypes.tex}
\end{table}

\begin{figure}[H]
\centering
\includegraphics[width=\linewidth]{generated/fig_archetypes.pdf}
\caption{Solved first-in mix of each archetype by seat.}
\label{fig:archetypes}
\end{figure}

Nine of the eleven archetypes contain at least two Hwang tiers at 10\% or more. ``Nut big cards''
is half Speculative and a third Marginal; ``Strong big pair'' is mostly Marginal; the low single-suited rundowns that Hwang calls
Premium fall into Rest.

\paragraph{Evaluation.} Table~\ref{tab:evaluation} compares the groupings in two ways. The first is the
EV each rule loses against the best action for each audited hand, with every group's action chosen
on one half of the hands and scored on the other half. The second is the share of variation in the
solver's fold rate, across all hands, that each grouping explains. The archetypes lose about a third
as much EV as Hwang's advice from under the gun to the cutoff and about half as much on the button,
and they come within a few hundredths of a big blind of the solver's own 780-bucket strategy. All
groupings do poorly in the small blind, where most hands are close to break-even.

\begin{table}[H]
\centering
\caption{How well each grouping reproduces the solved first-in decision.}
\label{tab:evaluation}
\small
\input{generated/table_evaluation.tex}
\end{table}

\section{Limitations}
\begin{itemize}
\item One stack depth (20 big blinds), pot-sized bets only, no rake, no antes. The archetypes are
fitted to this game and should be re-derived for deeper stacks, where implied odds favour Hwang's
straight hands.
\item Postflop play is abstracted and imperfect-recall, and multiplayer CFR has no equilibrium
guarantee. The only exploitability evidence is the small gain from changing the first action.
\item The audit follows the solved strategy after the first action, so it measures one-step
deviations, not full best responses.
\item The archetypes were designed after looking at the same solve that evaluates them. The
held-out split protects the choice of action per group but not the choice of rules; an independent
solve at another depth would be the stronger test.
\item Non-Premium tiers were audited on 400 hands each; many of those hands are too close to call at
2{,}048 deals per hand.
\end{itemize}

\section{Conclusion}
At 20 big blinds Hwang's tiers order hands correctly but misplace two families: low straight hands
that are not double-suited, which his scheme calls Premium or Speculative, and big-card hands with
a dangler, which it calls Marginal. The solver's first action is explained mainly by big cards, the
Ace suit and pair rank. Eleven archetypes built on those features reproduce the solver's first-in
choices far better than the four tiers, and nearly as well as the solver's own buckets, while
remaining simple enough to apply at the table.

\begin{thebibliography}{9}
\bibitem{hwang} J.~Hwang. \emph{Pot-Limit Omaha Poker: The Big Play Strategy}. Lyle Stuart
(Kensington), 2008, chapter 4.
\end{thebibliography}

\end{document}
